% 07_limitations -- written by Claude Code (D-045); revised after the peer review (Prompt 16: M2, M3h, M7, m1, m4, m15).
\section{Limitations}\label{sec:limitations}
\emph{Exam questions are a proxy.} MedQA items have one keyed answer and were written for a licensing exam; a tumor
board weighs imaging, pathology, staging, comorbidity, guidelines and the patient's wishes, often accepts more than one
option, and decides by discussion. Only \nM{} of our \nTest{} items concern management, and their result is imprecise.
Nothing here measures clinical usefulness or safety.

\emph{Contamination.} Most test items (\nOriginTrain{} of \nTest) come from MedQA's training split, which is widely
redistributed; whether the model saw them is unknown. The probe is weak (one hidden option, exact match, no control
set, many short options). Memorised items would raise accuracy and could make runs agree more, which would shrink, not
protect, the voting gain; similar gains in items from MedQA's train, dev and test splits argue against a large effect,
but those strata are small.

\emph{Screening by AI.} Oncology relevance was coded by an AI model and checked on a sample by a second, blind instance of
the same model, whose errors may be correlated with the first; the first screener adjudicated. The objective subset
gave a result in the same direction.

\emph{One model, prompt, option order and seed set.} The runs share a preference against option A, so their errors
are not independent; shuffling the options per run, a second seed set and a second model are the obvious next tests.

\emph{Statistics and literature.} Only RQ1 is confirmatory; other analyses are not corrected for multiple comparisons.
Fewer items were discordant than the power calculation assumed, so the estimate may be inflated, and the management
subgroup cannot detect gains below roughly 9--10 points. The 3-point threshold is our own choice. Several cited studies
were read only as abstracts, and we did not review the literature on benchmark contamination, selective prediction,
answer-position bias or GPU non-determinism.
